%%%%% PROJECT-SPECIFIC MACROS %%%%%
% Define your project-specific terms here
% Import with: %%%%% PROJECT-SPECIFIC MACROS %%%%%
% Define your project-specific terms here
% Import with: %%%%% PROJECT-SPECIFIC MACROS %%%%%
% Define your project-specific terms here
% Import with: %%%%% PROJECT-SPECIFIC MACROS %%%%%
% Define your project-specific terms here
% Import with: \input{commands/macros}
%
% INSTRUCTIONS:
% 1. Use \textsc{} for method/dataset/model names (small caps)
% 2. Add \xspace at the end for automatic spacing
% 3. Group related terms together
% 4. Use consistent naming conventions

\usepackage{xspace}

%%%%%%%%%%%%%%%%%%%%%%%%%%%%%%%%%%%%%%%%%%%%%%%%%%%%%%%%%%%%%%%%%
% METHOD NAMES
%%%%%%%%%%%%%%%%%%%%%%%%%%%%%%%%%%%%%%%%%%%%%%%%%%%%%%%%%%%%%%%%%

% Your method name (replace with actual name)
% Example: \newcommand{\ours}{{\bf YourMethod}\xspace}
% Example: \newcommand{\methodname}{\textsc{MethodName}\xspace}

%%%%%%%%%%%%%%%%%%%%%%%%%%%%%%%%%%%%%%%%%%%%%%%%%%%%%%%%%%%%%%%%%
% DATASET NAMES
%%%%%%%%%%%%%%%%%%%%%%%%%%%%%%%%%%%%%%%%%%%%%%%%%%%%%%%%%%%%%%%%%

% Dataset names (use \textsc for consistency)
% Example: \newcommand{\datasetone}{\textsc{Dataset-One}\xspace}
% Example: \newcommand{\datasettwo}{\textsc{Dataset-Two}\xspace}

%%%%%%%%%%%%%%%%%%%%%%%%%%%%%%%%%%%%%%%%%%%%%%%%%%%%%%%%%%%%%%%%%
% MODEL NAMES
%%%%%%%%%%%%%%%%%%%%%%%%%%%%%%%%%%%%%%%%%%%%%%%%%%%%%%%%%%%%%%%%%

% LLM and model names
% Example: \newcommand{\gpt}{\textsc{GPT-4o-mini}\xspace}
% Example: \newcommand{\llama}{\textsc{Llama-3.1-70B-Instruct}\xspace}
% Example: \newcommand{\claude}{\textsc{Claude-3.5-Sonnet}\xspace}

%%%%%%%%%%%%%%%%%%%%%%%%%%%%%%%%%%%%%%%%%%%%%%%%%%%%%%%%%%%%%%%%%
% BASELINE NAMES
%%%%%%%%%%%%%%%%%%%%%%%%%%%%%%%%%%%%%%%%%%%%%%%%%%%%%%%%%%%%%%%%%

% Baseline method names
% Example: \newcommand{\baseline}{\textsc{BaselineMethod}\xspace}
% Example: \newcommand{\fewshot}{\textsc{Few-shot}\xspace}

%%%%%%%%%%%%%%%%%%%%%%%%%%%%%%%%%%%%%%%%%%%%%%%%%%%%%%%%%%%%%%%%%
% ABBREVIATIONS
%%%%%%%%%%%%%%%%%%%%%%%%%%%%%%%%%%%%%%%%%%%%%%%%%%%%%%%%%%%%%%%%%

% Common abbreviations used in your paper
% Example: \newcommand{\ie}{i.e.,\xspace}
% Example: \newcommand{\eg}{e.g.,\xspace}
% Example: \newcommand{\cf}{cf.\xspace}
% Example: \newcommand{\wrt}{w.r.t.\xspace}

%%%%%%%%%%%%%%%%%%%%%%%%%%%%%%%%%%%%%%%%%%%%%%%%%%%%%%%%%%%%%%%%%
% DOMAIN-SPECIFIC TERMS
%%%%%%%%%%%%%%%%%%%%%%%%%%%%%%%%%%%%%%%%%%%%%%%%%%%%%%%%%%%%%%%%%

% Terms specific to your research domain
% Keep consistent formatting throughout the paper

%%%%%%%%%%%%%%%%%%%%%%%%%%%%%%%%%%%%%%%%%%%%%%%%%%%%%%%%%%%%%%%%%
% USAGE EXAMPLES
%%%%%%%%%%%%%%%%%%%%%%%%%%%%%%%%%%%%%%%%%%%%%%%%%%%%%%%%%%%%%%%%%
%
% GOOD:
%   \newcommand{\hypogenic}{\textsc{HypoGeniC}\xspace}
%   Usage: "We propose \hypogenic to generate hypotheses."
%   Result: "We propose HypoGeniC to generate hypotheses."
%
% GOOD:
%   \newcommand{\ours}{{\bf HypoGeniC}\xspace}
%   Usage: "Our method, \ours, outperforms baselines."
%   Result: "Our method, HypoGeniC, outperforms baselines."
%
% BAD (no \xspace):
%   \newcommand{\hypogenic}{\textsc{HypoGeniC}}
%   Usage: "\hypogenic outperforms..."
%   Result: "HypoGeniCoutperforms..." (missing space!)
%
%%%%%%%%%%%%%%%%%%%%%%%%%%%%%%%%%%%%%%%%%%%%%%%%%%%%%%%%%%%%%%%%%

%%%%%%%%%%%%%%%%%%%%%%%%%%%%%%%%%%%%%%%%%%%%%%%%%%%%%%%%%%%%%%%%%
% CARB PROJECT MACROS
%%%%%%%%%%%%%%%%%%%%%%%%%%%%%%%%%%%%%%%%%%%%%%%%%%%%%%%%%%%%%%%%%

% Benchmark
\newcommand{\carb}{\textsc{Carb}\xspace}

% Actions
\newcommand{\act}{\textsc{Act}\xspace}
\newcommand{\ask}{\textsc{Ask}\xspace}
\newcommand{\refuse}{\textsc{Refuse}\xspace}
\newcommand{\defer}{\textsc{Defer}\xspace}

% Elicitation regimes
\newcommand{\rzero}{\textsc{R0-Direct}\xspace}
\newcommand{\rone}{\textsc{R1-Afford}\xspace}
\newcommand{\rtwo}{\textsc{R2-Typed}\xspace}
\newcommand{\rthree}{\textsc{R3-Scalar}\xspace}
\newcommand{\rfour}{\textsc{R4-Recog}\xspace}

% Datasets
\newcommand{\coconot}{\textsc{CoCoNot}\xspace}
\newcommand{\clamber}{\textsc{Clamber}\xspace}
\newcommand{\inthree}{\textsc{IN3}\xspace}
\newcommand{\ambigqa}{\textsc{AmbigQA}\xspace}

% Models
\newcommand{\gptfive}{\textsc{GPT-5}\xspace}
\newcommand{\claude}{\textsc{Claude-Sonnet-4.5}\xspace}
\newcommand{\gemini}{\textsc{Gemini-2.5-Flash}\xspace}
\newcommand{\llama}{\textsc{Llama-3.3-70B}\xspace}
\newcommand{\qwen}{\textsc{Qwen3-4B}\xspace}
\newcommand{\qwenjudge}{\textsc{Qwen3-14B}\xspace}

% Abbreviations
\newcommand{\ie}{i.e.,\xspace}
\newcommand{\eg}{e.g.,\xspace}
\newcommand{\versus}{vs.\xspace}

%
% INSTRUCTIONS:
% 1. Use \textsc{} for method/dataset/model names (small caps)
% 2. Add \xspace at the end for automatic spacing
% 3. Group related terms together
% 4. Use consistent naming conventions

\usepackage{xspace}

%%%%%%%%%%%%%%%%%%%%%%%%%%%%%%%%%%%%%%%%%%%%%%%%%%%%%%%%%%%%%%%%%
% METHOD NAMES
%%%%%%%%%%%%%%%%%%%%%%%%%%%%%%%%%%%%%%%%%%%%%%%%%%%%%%%%%%%%%%%%%

% Your method name (replace with actual name)
% Example: \newcommand{\ours}{{\bf YourMethod}\xspace}
% Example: \newcommand{\methodname}{\textsc{MethodName}\xspace}

%%%%%%%%%%%%%%%%%%%%%%%%%%%%%%%%%%%%%%%%%%%%%%%%%%%%%%%%%%%%%%%%%
% DATASET NAMES
%%%%%%%%%%%%%%%%%%%%%%%%%%%%%%%%%%%%%%%%%%%%%%%%%%%%%%%%%%%%%%%%%

% Dataset names (use \textsc for consistency)
% Example: \newcommand{\datasetone}{\textsc{Dataset-One}\xspace}
% Example: \newcommand{\datasettwo}{\textsc{Dataset-Two}\xspace}

%%%%%%%%%%%%%%%%%%%%%%%%%%%%%%%%%%%%%%%%%%%%%%%%%%%%%%%%%%%%%%%%%
% MODEL NAMES
%%%%%%%%%%%%%%%%%%%%%%%%%%%%%%%%%%%%%%%%%%%%%%%%%%%%%%%%%%%%%%%%%

% LLM and model names
% Example: \newcommand{\gpt}{\textsc{GPT-4o-mini}\xspace}
% Example: \newcommand{\llama}{\textsc{Llama-3.1-70B-Instruct}\xspace}
% Example: \newcommand{\claude}{\textsc{Claude-3.5-Sonnet}\xspace}

%%%%%%%%%%%%%%%%%%%%%%%%%%%%%%%%%%%%%%%%%%%%%%%%%%%%%%%%%%%%%%%%%
% BASELINE NAMES
%%%%%%%%%%%%%%%%%%%%%%%%%%%%%%%%%%%%%%%%%%%%%%%%%%%%%%%%%%%%%%%%%

% Baseline method names
% Example: \newcommand{\baseline}{\textsc{BaselineMethod}\xspace}
% Example: \newcommand{\fewshot}{\textsc{Few-shot}\xspace}

%%%%%%%%%%%%%%%%%%%%%%%%%%%%%%%%%%%%%%%%%%%%%%%%%%%%%%%%%%%%%%%%%
% ABBREVIATIONS
%%%%%%%%%%%%%%%%%%%%%%%%%%%%%%%%%%%%%%%%%%%%%%%%%%%%%%%%%%%%%%%%%

% Common abbreviations used in your paper
% Example: \newcommand{\ie}{i.e.,\xspace}
% Example: \newcommand{\eg}{e.g.,\xspace}
% Example: \newcommand{\cf}{cf.\xspace}
% Example: \newcommand{\wrt}{w.r.t.\xspace}

%%%%%%%%%%%%%%%%%%%%%%%%%%%%%%%%%%%%%%%%%%%%%%%%%%%%%%%%%%%%%%%%%
% DOMAIN-SPECIFIC TERMS
%%%%%%%%%%%%%%%%%%%%%%%%%%%%%%%%%%%%%%%%%%%%%%%%%%%%%%%%%%%%%%%%%

% Terms specific to your research domain
% Keep consistent formatting throughout the paper

%%%%%%%%%%%%%%%%%%%%%%%%%%%%%%%%%%%%%%%%%%%%%%%%%%%%%%%%%%%%%%%%%
% USAGE EXAMPLES
%%%%%%%%%%%%%%%%%%%%%%%%%%%%%%%%%%%%%%%%%%%%%%%%%%%%%%%%%%%%%%%%%
%
% GOOD:
%   \newcommand{\hypogenic}{\textsc{HypoGeniC}\xspace}
%   Usage: "We propose \hypogenic to generate hypotheses."
%   Result: "We propose HypoGeniC to generate hypotheses."
%
% GOOD:
%   \newcommand{\ours}{{\bf HypoGeniC}\xspace}
%   Usage: "Our method, \ours, outperforms baselines."
%   Result: "Our method, HypoGeniC, outperforms baselines."
%
% BAD (no \xspace):
%   \newcommand{\hypogenic}{\textsc{HypoGeniC}}
%   Usage: "\hypogenic outperforms..."
%   Result: "HypoGeniCoutperforms..." (missing space!)
%
%%%%%%%%%%%%%%%%%%%%%%%%%%%%%%%%%%%%%%%%%%%%%%%%%%%%%%%%%%%%%%%%%

%%%%%%%%%%%%%%%%%%%%%%%%%%%%%%%%%%%%%%%%%%%%%%%%%%%%%%%%%%%%%%%%%
% CARB PROJECT MACROS
%%%%%%%%%%%%%%%%%%%%%%%%%%%%%%%%%%%%%%%%%%%%%%%%%%%%%%%%%%%%%%%%%

% Benchmark
\newcommand{\carb}{\textsc{Carb}\xspace}

% Actions
\newcommand{\act}{\textsc{Act}\xspace}
\newcommand{\ask}{\textsc{Ask}\xspace}
\newcommand{\refuse}{\textsc{Refuse}\xspace}
\newcommand{\defer}{\textsc{Defer}\xspace}

% Elicitation regimes
\newcommand{\rzero}{\textsc{R0-Direct}\xspace}
\newcommand{\rone}{\textsc{R1-Afford}\xspace}
\newcommand{\rtwo}{\textsc{R2-Typed}\xspace}
\newcommand{\rthree}{\textsc{R3-Scalar}\xspace}
\newcommand{\rfour}{\textsc{R4-Recog}\xspace}

% Datasets
\newcommand{\coconot}{\textsc{CoCoNot}\xspace}
\newcommand{\clamber}{\textsc{Clamber}\xspace}
\newcommand{\inthree}{\textsc{IN3}\xspace}
\newcommand{\ambigqa}{\textsc{AmbigQA}\xspace}

% Models
\newcommand{\gptfive}{\textsc{GPT-5}\xspace}
\newcommand{\claude}{\textsc{Claude-Sonnet-4.5}\xspace}
\newcommand{\gemini}{\textsc{Gemini-2.5-Flash}\xspace}
\newcommand{\llama}{\textsc{Llama-3.3-70B}\xspace}
\newcommand{\qwen}{\textsc{Qwen3-4B}\xspace}
\newcommand{\qwenjudge}{\textsc{Qwen3-14B}\xspace}

% Abbreviations
\newcommand{\ie}{i.e.,\xspace}
\newcommand{\eg}{e.g.,\xspace}
\newcommand{\versus}{vs.\xspace}

%
% INSTRUCTIONS:
% 1. Use \textsc{} for method/dataset/model names (small caps)
% 2. Add \xspace at the end for automatic spacing
% 3. Group related terms together
% 4. Use consistent naming conventions

\usepackage{xspace}

%%%%%%%%%%%%%%%%%%%%%%%%%%%%%%%%%%%%%%%%%%%%%%%%%%%%%%%%%%%%%%%%%
% METHOD NAMES
%%%%%%%%%%%%%%%%%%%%%%%%%%%%%%%%%%%%%%%%%%%%%%%%%%%%%%%%%%%%%%%%%

% Your method name (replace with actual name)
% Example: \newcommand{\ours}{{\bf YourMethod}\xspace}
% Example: \newcommand{\methodname}{\textsc{MethodName}\xspace}

%%%%%%%%%%%%%%%%%%%%%%%%%%%%%%%%%%%%%%%%%%%%%%%%%%%%%%%%%%%%%%%%%
% DATASET NAMES
%%%%%%%%%%%%%%%%%%%%%%%%%%%%%%%%%%%%%%%%%%%%%%%%%%%%%%%%%%%%%%%%%

% Dataset names (use \textsc for consistency)
% Example: \newcommand{\datasetone}{\textsc{Dataset-One}\xspace}
% Example: \newcommand{\datasettwo}{\textsc{Dataset-Two}\xspace}

%%%%%%%%%%%%%%%%%%%%%%%%%%%%%%%%%%%%%%%%%%%%%%%%%%%%%%%%%%%%%%%%%
% MODEL NAMES
%%%%%%%%%%%%%%%%%%%%%%%%%%%%%%%%%%%%%%%%%%%%%%%%%%%%%%%%%%%%%%%%%

% LLM and model names
% Example: \newcommand{\gpt}{\textsc{GPT-4o-mini}\xspace}
% Example: \newcommand{\llama}{\textsc{Llama-3.1-70B-Instruct}\xspace}
% Example: \newcommand{\claude}{\textsc{Claude-3.5-Sonnet}\xspace}

%%%%%%%%%%%%%%%%%%%%%%%%%%%%%%%%%%%%%%%%%%%%%%%%%%%%%%%%%%%%%%%%%
% BASELINE NAMES
%%%%%%%%%%%%%%%%%%%%%%%%%%%%%%%%%%%%%%%%%%%%%%%%%%%%%%%%%%%%%%%%%

% Baseline method names
% Example: \newcommand{\baseline}{\textsc{BaselineMethod}\xspace}
% Example: \newcommand{\fewshot}{\textsc{Few-shot}\xspace}

%%%%%%%%%%%%%%%%%%%%%%%%%%%%%%%%%%%%%%%%%%%%%%%%%%%%%%%%%%%%%%%%%
% ABBREVIATIONS
%%%%%%%%%%%%%%%%%%%%%%%%%%%%%%%%%%%%%%%%%%%%%%%%%%%%%%%%%%%%%%%%%

% Common abbreviations used in your paper
% Example: \newcommand{\ie}{i.e.,\xspace}
% Example: \newcommand{\eg}{e.g.,\xspace}
% Example: \newcommand{\cf}{cf.\xspace}
% Example: \newcommand{\wrt}{w.r.t.\xspace}

%%%%%%%%%%%%%%%%%%%%%%%%%%%%%%%%%%%%%%%%%%%%%%%%%%%%%%%%%%%%%%%%%
% DOMAIN-SPECIFIC TERMS
%%%%%%%%%%%%%%%%%%%%%%%%%%%%%%%%%%%%%%%%%%%%%%%%%%%%%%%%%%%%%%%%%

% Terms specific to your research domain
% Keep consistent formatting throughout the paper

%%%%%%%%%%%%%%%%%%%%%%%%%%%%%%%%%%%%%%%%%%%%%%%%%%%%%%%%%%%%%%%%%
% USAGE EXAMPLES
%%%%%%%%%%%%%%%%%%%%%%%%%%%%%%%%%%%%%%%%%%%%%%%%%%%%%%%%%%%%%%%%%
%
% GOOD:
%   \newcommand{\hypogenic}{\textsc{HypoGeniC}\xspace}
%   Usage: "We propose \hypogenic to generate hypotheses."
%   Result: "We propose HypoGeniC to generate hypotheses."
%
% GOOD:
%   \newcommand{\ours}{{\bf HypoGeniC}\xspace}
%   Usage: "Our method, \ours, outperforms baselines."
%   Result: "Our method, HypoGeniC, outperforms baselines."
%
% BAD (no \xspace):
%   \newcommand{\hypogenic}{\textsc{HypoGeniC}}
%   Usage: "\hypogenic outperforms..."
%   Result: "HypoGeniCoutperforms..." (missing space!)
%
%%%%%%%%%%%%%%%%%%%%%%%%%%%%%%%%%%%%%%%%%%%%%%%%%%%%%%%%%%%%%%%%%

%%%%%%%%%%%%%%%%%%%%%%%%%%%%%%%%%%%%%%%%%%%%%%%%%%%%%%%%%%%%%%%%%
% CARB PROJECT MACROS
%%%%%%%%%%%%%%%%%%%%%%%%%%%%%%%%%%%%%%%%%%%%%%%%%%%%%%%%%%%%%%%%%

% Benchmark
\newcommand{\carb}{\textsc{Carb}\xspace}

% Actions
\newcommand{\act}{\textsc{Act}\xspace}
\newcommand{\ask}{\textsc{Ask}\xspace}
\newcommand{\refuse}{\textsc{Refuse}\xspace}
\newcommand{\defer}{\textsc{Defer}\xspace}

% Elicitation regimes
\newcommand{\rzero}{\textsc{R0-Direct}\xspace}
\newcommand{\rone}{\textsc{R1-Afford}\xspace}
\newcommand{\rtwo}{\textsc{R2-Typed}\xspace}
\newcommand{\rthree}{\textsc{R3-Scalar}\xspace}
\newcommand{\rfour}{\textsc{R4-Recog}\xspace}

% Datasets
\newcommand{\coconot}{\textsc{CoCoNot}\xspace}
\newcommand{\clamber}{\textsc{Clamber}\xspace}
\newcommand{\inthree}{\textsc{IN3}\xspace}
\newcommand{\ambigqa}{\textsc{AmbigQA}\xspace}

% Models
\newcommand{\gptfive}{\textsc{GPT-5}\xspace}
\newcommand{\claude}{\textsc{Claude-Sonnet-4.5}\xspace}
\newcommand{\gemini}{\textsc{Gemini-2.5-Flash}\xspace}
\newcommand{\llama}{\textsc{Llama-3.3-70B}\xspace}
\newcommand{\qwen}{\textsc{Qwen3-4B}\xspace}
\newcommand{\qwenjudge}{\textsc{Qwen3-14B}\xspace}

% Abbreviations
\newcommand{\ie}{i.e.,\xspace}
\newcommand{\eg}{e.g.,\xspace}
\newcommand{\versus}{vs.\xspace}

%
% INSTRUCTIONS:
% 1. Use \textsc{} for method/dataset/model names (small caps)
% 2. Add \xspace at the end for automatic spacing
% 3. Group related terms together
% 4. Use consistent naming conventions

\usepackage{xspace}

%%%%%%%%%%%%%%%%%%%%%%%%%%%%%%%%%%%%%%%%%%%%%%%%%%%%%%%%%%%%%%%%%
% METHOD NAMES
%%%%%%%%%%%%%%%%%%%%%%%%%%%%%%%%%%%%%%%%%%%%%%%%%%%%%%%%%%%%%%%%%

% Your method name (replace with actual name)
% Example: \newcommand{\ours}{{\bf YourMethod}\xspace}
% Example: \newcommand{\methodname}{\textsc{MethodName}\xspace}

%%%%%%%%%%%%%%%%%%%%%%%%%%%%%%%%%%%%%%%%%%%%%%%%%%%%%%%%%%%%%%%%%
% DATASET NAMES
%%%%%%%%%%%%%%%%%%%%%%%%%%%%%%%%%%%%%%%%%%%%%%%%%%%%%%%%%%%%%%%%%

% Dataset names (use \textsc for consistency)
% Example: \newcommand{\datasetone}{\textsc{Dataset-One}\xspace}
% Example: \newcommand{\datasettwo}{\textsc{Dataset-Two}\xspace}

%%%%%%%%%%%%%%%%%%%%%%%%%%%%%%%%%%%%%%%%%%%%%%%%%%%%%%%%%%%%%%%%%
% MODEL NAMES
%%%%%%%%%%%%%%%%%%%%%%%%%%%%%%%%%%%%%%%%%%%%%%%%%%%%%%%%%%%%%%%%%

% LLM and model names
% Example: \newcommand{\gpt}{\textsc{GPT-4o-mini}\xspace}
% Example: \newcommand{\llama}{\textsc{Llama-3.1-70B-Instruct}\xspace}
% Example: \newcommand{\claude}{\textsc{Claude-3.5-Sonnet}\xspace}

%%%%%%%%%%%%%%%%%%%%%%%%%%%%%%%%%%%%%%%%%%%%%%%%%%%%%%%%%%%%%%%%%
% BASELINE NAMES
%%%%%%%%%%%%%%%%%%%%%%%%%%%%%%%%%%%%%%%%%%%%%%%%%%%%%%%%%%%%%%%%%

% Baseline method names
% Example: \newcommand{\baseline}{\textsc{BaselineMethod}\xspace}
% Example: \newcommand{\fewshot}{\textsc{Few-shot}\xspace}

%%%%%%%%%%%%%%%%%%%%%%%%%%%%%%%%%%%%%%%%%%%%%%%%%%%%%%%%%%%%%%%%%
% ABBREVIATIONS
%%%%%%%%%%%%%%%%%%%%%%%%%%%%%%%%%%%%%%%%%%%%%%%%%%%%%%%%%%%%%%%%%

% Common abbreviations used in your paper
% Example: \newcommand{\ie}{i.e.,\xspace}
% Example: \newcommand{\eg}{e.g.,\xspace}
% Example: \newcommand{\cf}{cf.\xspace}
% Example: \newcommand{\wrt}{w.r.t.\xspace}

%%%%%%%%%%%%%%%%%%%%%%%%%%%%%%%%%%%%%%%%%%%%%%%%%%%%%%%%%%%%%%%%%
% DOMAIN-SPECIFIC TERMS
%%%%%%%%%%%%%%%%%%%%%%%%%%%%%%%%%%%%%%%%%%%%%%%%%%%%%%%%%%%%%%%%%

% Terms specific to your research domain
% Keep consistent formatting throughout the paper

%%%%%%%%%%%%%%%%%%%%%%%%%%%%%%%%%%%%%%%%%%%%%%%%%%%%%%%%%%%%%%%%%
% USAGE EXAMPLES
%%%%%%%%%%%%%%%%%%%%%%%%%%%%%%%%%%%%%%%%%%%%%%%%%%%%%%%%%%%%%%%%%
%
% GOOD:
%   \newcommand{\hypogenic}{\textsc{HypoGeniC}\xspace}
%   Usage: "We propose \hypogenic to generate hypotheses."
%   Result: "We propose HypoGeniC to generate hypotheses."
%
% GOOD:
%   \newcommand{\ours}{{\bf HypoGeniC}\xspace}
%   Usage: "Our method, \ours, outperforms baselines."
%   Result: "Our method, HypoGeniC, outperforms baselines."
%
% BAD (no \xspace):
%   \newcommand{\hypogenic}{\textsc{HypoGeniC}}
%   Usage: "\hypogenic outperforms..."
%   Result: "HypoGeniCoutperforms..." (missing space!)
%
%%%%%%%%%%%%%%%%%%%%%%%%%%%%%%%%%%%%%%%%%%%%%%%%%%%%%%%%%%%%%%%%%

%%%%%%%%%%%%%%%%%%%%%%%%%%%%%%%%%%%%%%%%%%%%%%%%%%%%%%%%%%%%%%%%%
% CARB PROJECT MACROS
%%%%%%%%%%%%%%%%%%%%%%%%%%%%%%%%%%%%%%%%%%%%%%%%%%%%%%%%%%%%%%%%%

% Benchmark
\newcommand{\carb}{\textsc{Carb}\xspace}

% Actions
\newcommand{\act}{\textsc{Act}\xspace}
\newcommand{\ask}{\textsc{Ask}\xspace}
\newcommand{\refuse}{\textsc{Refuse}\xspace}
\newcommand{\defer}{\textsc{Defer}\xspace}

% Elicitation regimes
\newcommand{\rzero}{\textsc{R0-Direct}\xspace}
\newcommand{\rone}{\textsc{R1-Afford}\xspace}
\newcommand{\rtwo}{\textsc{R2-Typed}\xspace}
\newcommand{\rthree}{\textsc{R3-Scalar}\xspace}
\newcommand{\rfour}{\textsc{R4-Recog}\xspace}

% Datasets
\newcommand{\coconot}{\textsc{CoCoNot}\xspace}
\newcommand{\clamber}{\textsc{Clamber}\xspace}
\newcommand{\inthree}{\textsc{IN3}\xspace}
\newcommand{\ambigqa}{\textsc{AmbigQA}\xspace}

% Models
\newcommand{\gptfive}{\textsc{GPT-5}\xspace}
\newcommand{\claude}{\textsc{Claude-Sonnet-4.5}\xspace}
\newcommand{\gemini}{\textsc{Gemini-2.5-Flash}\xspace}
\newcommand{\llama}{\textsc{Llama-3.3-70B}\xspace}
\newcommand{\qwen}{\textsc{Qwen3-4B}\xspace}
\newcommand{\qwenjudge}{\textsc{Qwen3-14B}\xspace}

% Abbreviations
\newcommand{\ie}{i.e.,\xspace}
\newcommand{\eg}{e.g.,\xspace}
\newcommand{\versus}{vs.\xspace}
